\documentclass[10pt,a4paper]{amsart}
\usepackage[margin=19mm]{geometry}
\usepackage{amsmath,amssymb,mathtools}
\usepackage[scr=rsfs,cal=euler]{mathalfa}
\usepackage{xfrac}
\usepackage[unicode,hidelinks,bookmarks=false]{hyperref}
\newcommand{\ie}{i.e.\ }
\newcommand{\cf}{cf.\ }
\newcommand{\resp}{respectively\ }
\newcommand{\Subsection}{\subsection}
\providecommand{\nobreakdash}{\nobreak\hbox{-}}
\setlength{\emergencystretch}{2em}
\multlinegap0pt
\usepackage{bm}
\usepackage{bbm}
\usepackage{dsfont}
\usepackage{xspace}
\usepackage{cancel}
\usepackage{mathtools}
\usepackage{amsthm}
\makeatletter
\@ifundefined{theorem}{\newtheorem{theorem}{Theorem}}{}
\@ifundefined{lemma}{\newtheorem{lemma}[theorem]{Lemma}}{}
\makeatother
\newcommand\ZZ{\mathbb{Z}}
\title[Pairs of intertwined integer sequences]{Pairs of intertwined integer sequences}
\author{Christian Kassel, Christophe Reutenauer}
\date{Source: arXiv:2507.15780v2 (CC BY 4.0)}
\begin{document}
\maketitle

\noindent\emph{Source and licence.} Pairs of intertwined integer sequences. Version arXiv:2507.15780v2 (CC BY 4.0), distributed under the Creative Commons Attribution 4.0 International licence, as recorded in the arXiv licence metadata for this exact version and shown on the abstract page https://arxiv.org/abs/2507.15780v2. Full rights notice: licenses/arxiv-2507.15780-CC-BY-4.0.txt.\par\smallskip

\noindent\textit{Editorial scope.} Counted unit: the Lemma whose source label is \texttt{formula1} in the article (source0.tex lines 1129--1136, source label \texttt{formula1}), delivered together with the article's own complete original proof of it (source0.tex lines 1138--1154, one \texttt{proof} environment of 17 lines). No mathematics is written, repaired or re-derived: statement and proof are the article's own text line for line. Required context retained uncounted: formula-cn0 (lines 1100--1102); formula-cni (lines 1105--1116). Every definition, display and label of those lines is retained unchanged. Omitted from this package and not redistributed: 1--1099, 1103--1104, 1117--1128, 1137--1137, 1155--1402. Nothing inside a retained excerpt is omitted or altered: every retained body is delivered exactly as the article's lines are written, so no omission receipt of any kind accompanies this package. Citations: the retained excerpts carry no citation command at all, so no bibliography entry of the article is transcribed and no reference is redistributed. Reference adaptation: none was needed and none was made; every reference inside the delivered unit resolves inside it, so no unresolved reference or citation is produced. Printed slips kept literal: no printed word, symbol, hypothesis or number of the counted statement or of its proof was altered. Layout and class: the article's own class is \texttt{amsart} with packages including amsmath, amsthm, amssymb, txfonts, mathabx, hyperref; the wrapper is the lane's pinned \texttt{amsart} editorial wrapper at 10pt on A4, which restates the theorem-like environments the retained text needs and adds the article's own macro definitions that the retained text needs (\texttt{\textbackslash ZZ}), with extra wrapper packages (bm, bbm, dsfont, xspace, cancel, mathtools). The delivered numbering is the unit's own. Assets: no retained span contains a figure environment, a graphics-inclusion command, a PDF-page inclusion, a table or a TikZ picture; no image file is copied into this package, the package declares no asset dependency and no image file is redistributed with it. Licence: version arXiv:2507.15780v2 of this article is distributed under the Creative Commons Attribution 4.0 International licence, as recorded in the arXiv licence metadata for this exact version and shown on the abstract page https://arxiv.org/abs/2507.15780v2. A full rights notice is delivered at licenses/arxiv-2507.15780-CC-BY-4.0.txt. Characters: every retained excerpt is pure ASCII, so the delivered engine, XeLaTeX with its default Latin Modern text font, reports no missing character.
\begin{equation}\label{formula-cn0}
c_{n,0} = 2 (-1)^r
\end{equation}

\begin{equation}\label{formula-cni}
c_{n,i} =
\left\{
\begin{array}{cl}
(-1)^k & \text{if}\;\, n = k(k+2i +1)/2 \;\; \text{for some integer}\; k  \geq 1, \\
\noalign{\smallskip}
(-1)^{k-1} & \text{if}\;\, n = k(k+2i -1)/2 \;\; \text{for some integer}\; k \geq 1, \\
\noalign{\smallskip}
0 & \text{otherwise.} 
\end{array}
\right.
\end{equation}

\begin{lemma}\label{formula1} 
We have
\[
\frac{C_n(q)}{q^n} = \sum\,  (-1)^h \, \left( q^a + q^{-a} \right), 
\]
where the summation is over all $h\in\ZZ_{\geq 1}$ and $a\in \ZZ$ such that $n=h(h+1+2a)/2$. 
Moreover, if $n$ has such a representation and $|a|$ is given, then $a,h$ are unique.
\end{lemma}

\begin{proof}
From~\eqref{formula-cn0} it follows that both sides have the same constant term. 

From now on we concentrate on the case when $a$ is a nonzero integer.
Let $i\geq 1$.
Suppose that $n=k(k+2i +1)/2$ with $k\geq 1$ as in the first row of~\eqref{formula-cni}; 
then $n=h(h+2a +1)/2$ with $h=k \geq 1$, $a=i\in \ZZ \setminus \{0\}$, and $(-1)^k=(-1)^h$. 

Now suppose that $n = k(k+2i -1)/2$ with $k\geq 1$ as in the second row of~\eqref{formula-cni}; 
then $n=h(h+2a +1)/2$ with $h=k+2i-1\geq 1$, $a=-i\in \ZZ \setminus \{0\}$, and $(-1)^{k-1}=(-1)^h$. 

Conversely, let $n=h(h+2a +1)/2$ with $h\geq 1$, $i\in \ZZ$ and $a\neq 0$. 
Then either $a\geq 1$, and then $n=k(k+2i +1)/2$ with $k=h$, $i=a\geq 1$, and $(-1)^k=(-1)^h$; 
or $a < 0$, and then $n=k(k+2i-1)/2$, $k=h+2a+1\geq 1$, $i = -a > 0$, and $(-1)^{k-1}=(-1)^h$.

This proves the lemma in view of the observation made just before.
\end{proof}

\end{document}
